\scanpage{228}

Parts (i) and (ii) yield the

\uline{Corollary 1} The following conditions are equivalent:

(i) There exists an open $U\quad y$ such that $X\ U$ is projective over $U$

(ii) For every $n$, $X_n$ is a projective analytic space.

For instance, if $\dim X_0\leq 1$, then (ii) and hence (i) holds.

The proof of the theorem only uses Grauerts analogues of the algebraic theorems of finiteness and comparison for direct images (cf his blue paper) and the usual exact sequence $0\quad\underline{Z}\quad\underline{O}\quad\underline{O}^{\bullet}\quad0$, together with some standard use of Mittag-Leffler story and five lemma. It is valid in fact for any $H^i(\underline{O}^{\bullet})$, not only $i=1$ (which seems the only one however to have geometric significance).

In the case of a \uline{formal} scheme propr over a complete noeth. local ring with residue field of \uline{caracteristic} 0, the analogon of the previous theorem (reducing to statements (i), (i bis)) \deleted{\unsure{(iii))|(iii)}} hold true, and I wrote a purely algebraic proof of this, relying only on the fact that the Kernel and Cokernel of $\underline{\operatorname{Pic}}_{X_{n+1}}\qquad\underline{\operatorname{Pic}}_{X_n}$ are without torsion, and Néron's finiteness theorem; in particular, the analogon of corollary 1 holds true in this case. These results break down of course in car. $>0$.

However, using the (as yet unwritten !) Gaga of Serre-Grauert-Remmert-Grothendieck (cf Grauert-Remmerts paper, complemented by the method of an old talk of mine in Cartan's Seminar, to recover the case of proper morphisms of \unsure{shhemes} from the projective one, via Chow's lemma \ldots), the analytic theorem above yields an interesting intrinsic property of analytic algebras over $\underline{C}$, with respect to algebraic geometry

\margin{227}
